\begin{center}
\SourceLeaf{029}\pdfbookmark[2]{11.}{article-011}11.
\end{center}
Problema inversum, celebre sub nomine \emph{problematis} \textsc{Kepler}\emph{i}, scilicet ex anomalia media invenire veram atque radium vectorem, longe frequentioris usus est. Astronomi aequationem centri per seriem infinitam secundum sinus angulorum $M$, $2M$, $3M$ etc. progredientem exhibere solent, quorum sinuum coëfficientes singuli et ipsi sunt series secundum potestates excentricitatis in infinitum excurrentes. Huic formulae pro aequatione centri, quam plures auctores evolverunt, hic immorari eo minus necessarium duximus, quod, nostro quidem iudicio, ad usum practicum, praesertim si excentricitas perparva non fuerit, longe minus idonea est, quam methodus indirecta, quam itaque in ea forma, quae maxime commoda nobis videtur, aliquanto fusius explicabimus.

Aequatio XII, $E = M + e\sin E$, quae ad transcendentium genus referenda est solutionemque per operationes finitas directas non admittit, tentando solvenda est, incipiendo a valore quodam approximato ipsius $E$, qui per methodos idoneas toties repetitas corrigitur, usque dum illi aequationi exacte satisfaciat, i.\,e. vel omni quam tabulae sinuum permittunt praecisione, vel ea saltem, quae ad scopum propositum sufficit. Quodsi hae correctiones haud temere sed per normam tutam atque certam instituuntur, vix ullum discrimen essentiale inter methodum talem indirectam atque solutionem per series adest, nisi quod in illa valor primus incognitae aliquatenus est arbitrarius, quod potius pro lucro habendum, quum valor apte electus correctiones insigniter accelerare permittat. Supponamus, $\epsilon$ esse valorem approximatum ipsius $E$, atque $x$ correctionem illi adhuc adiiciendam (in secundis expressam), ita ut valor $E = \epsilon + x$ aequationi nostrae exacte satisfaciat. Computetur $e\sin\epsilon$ in secundis per logarithmos, quod dum perficitur, simul e tabulis notetur variatio ipsius $\log\sin\epsilon$ pro $1^{\prime\prime}$ variatione ipsius $\epsilon$, atque variatio $\log e\sin\epsilon$ pro variatione unius unitatis in numero $e\sin\epsilon$; sint hae variationes sine respectu signorum resp. $\lambda$, $\mu$, ubi vix opus est monere, utrumque logarithmum per aeque multas figuras decimales expressum supponi. Quodsi iam $\epsilon$ ad verum ipsius $E$ valorem tam prope iam accedit, ut variationes logarithmi sinus ab $\epsilon$ usque ad $\epsilon+x$, variationesque logarithmi numeri ab $e\sin\epsilon$ usque ad $e\sin(\epsilon+x)$ pro uniformibus habere liceat, manifesto statui poterit $e\sin(\epsilon+x) = e\sin\epsilon \pm \frac{\lambda x}{\mu}$, signo \SourceLeaf{030}superiori pro quadrante primo et quarto, inferiori pro secundo et tertio valente. Quare quum sit $\epsilon+x = M+e\sin(\epsilon+x)$, fit $x = \frac{\mu}{\mu\mp\lambda}(M+e\sin\epsilon-\epsilon)$, valorque verus ipsius $E$ sive $\epsilon+x = M+e\sin\epsilon \pm \frac{\lambda}{\mu\mp\lambda}(M+e\sin\epsilon-\epsilon)$, signis ea qua diximus ratione determinatis. Ceterum facile perspicitur, esse sine respectu signi $\mu:\lambda = 1:e\cos\epsilon$, adeoque semper $\mu > \lambda$, unde concluditur, in quadrante primo et ultimo $M+e\sin\epsilon$ iacere inter $\epsilon$ atque $\epsilon+x$, in secundo ac tertio vero $\epsilon+x$ inter $\epsilon$ atque $M+e\sin\epsilon$, quae regula attentionem ad signa sublevare potest. Si valor suppositus $\epsilon$ nimis adhuc a vero aberraverat, quam ut suppositionem supra traditam pro satis exacta habere liceret, certe per hanc methodum invenietur valor multo propior, quo eadem operatio iterum adhuc, pluriesve si opus videtur, repetenda erit. Nullo vero negotio patet, si differentia valoris primi $\epsilon$ a vero tamquam quantitas ordinis primi spectetur, errorem valoris novi ad ordinem secundum referendum fore, et per operationem iteratam ad ordinem quartum, octavum etc. deprimi. Quo minor insuper fuerit excentricitas, eo velocius correctiones successivae convergent.
